% Options for packages loaded elsewhere
\PassOptionsToPackage{unicode}{hyperref}
\PassOptionsToPackage{hyphens}{url}
\documentclass[
]{article}
\usepackage{xcolor}
\usepackage{amsmath,amssymb}
\setcounter{secnumdepth}{-\maxdimen} % remove section numbering
\usepackage{iftex}
\ifPDFTeX
  \usepackage[T1]{fontenc}
  \usepackage[utf8]{inputenc}
  \usepackage{textcomp} % provide euro and other symbols
\else % if luatex or xetex
  \usepackage{unicode-math} % this also loads fontspec
  \defaultfontfeatures{Scale=MatchLowercase}
  \defaultfontfeatures[\rmfamily]{Ligatures=TeX,Scale=1}
\fi
\usepackage{lmodern}
\ifPDFTeX\else
  % xetex/luatex font selection
\fi
% Use upquote if available, for straight quotes in verbatim environments
\IfFileExists{upquote.sty}{\usepackage{upquote}}{}
\IfFileExists{microtype.sty}{% use microtype if available
  \usepackage[]{microtype}
  \UseMicrotypeSet[protrusion]{basicmath} % disable protrusion for tt fonts
}{}
\makeatletter
\@ifundefined{KOMAClassName}{% if non-KOMA class
  \IfFileExists{parskip.sty}{%
    \usepackage{parskip}
  }{% else
    \setlength{\parindent}{0pt}
    \setlength{\parskip}{6pt plus 2pt minus 1pt}}
}{% if KOMA class
  \KOMAoptions{parskip=half}}
\makeatother
\setlength{\emergencystretch}{3em} % prevent overfull lines
\providecommand{\tightlist}{%
  \setlength{\itemsep}{0pt}\setlength{\parskip}{0pt}}
\usepackage{bookmark}
\IfFileExists{xurl.sty}{\usepackage{xurl}}{} % add URL line breaks if available
\urlstyle{same}
\hypersetup{
  hidelinks,
  pdfcreator={LaTeX via pandoc}}

\author{}
\date{}

\begin{document}

\textbf{Abstract}

\begin{quote}
Many Kabaddi players, especially those who are still learning the
basics, do not always train under a coach. Most of the time, they
practice on their own or with limited guidance, which makes it hard to
tell whether their movements are actually correct. Even small mistakes
in posture or timing can go unnoticed and affect progress over time. To
help learners during such practice sessions, the project uses an
AR-based Kabaddi Ghost Trainer that displays correct movements visually
and provides automated feedback for improvement.

The proposed system begins by extracting body poses from recorded videos
of expert Kabaddi players using deep learning--based pose estimation
methods. These poses are then cleaned and organized to create a
reference movement sequence, referred to as a ``ghost'' motion. During
training, this ghost motion is shown in an augmented reality environment
so that learners can clearly see and follow the expert's movement while
practicing. At the same time, the system continuously captures the
learner's own body movements for analysis.

The learner's motion is compared with the reference sequence by
observing joint alignment, posture consistency, and movement timing
across frames. Rather than presenting only raw scores, the system
converts these comparisons into simple and meaningful feedback using an
LLM-based module. This feedback helps learners understand their mistakes
and make improvements.

In summary, the project brings together pose estimation, augmented
reality, and feedback generation to support individual Kabaddi practice,
showing how such technologies can be used to assist training in sports
that depend heavily on movement and technique.
\end{quote}

\textbf{Chapter 1}

\textbf{INTRODUCTION}

\textbf{1.1 OVERVIEW}

\begin{quote}
Kabaddi is a physically demanding sport that requires precise
coordination between body posture, timing, balance, and spatial
awareness. Effective execution of techniques such as raids, touches, and
evasive movements depends not only on physical strength but also on
accurate replication of expert movement patterns. Traditional training
methods largely rely on in-person coaching, verbal instructions, and
repeated observation of demonstrations, which may not always provide
objective or measurable feedback on performance quality. In situations
where direct supervision is limited, learners often struggle to identify
subtle posture deviations or timing errors during practice sessions.

This project addresses the need for an intelligent training support
system by integrating computer vision, pose analysis, and augmented
reality technologies. The core idea is to provide a structured and
measurable training environment where expert movements can be visualized
and compared against user-performed actions. By focusing on pose-level
analysis rather than raw video comparison, this work emphasizes
explainability and consistency in performance evaluation. The system is
designed to function as a virtual training aid that assists learners in
understanding movement accuracy, identifying deviations, and improving
technique execution in a systematic and data-driven manner.
\end{quote}

\textbf{1.2 PROBLEM STATEMENT}

\begin{quote}
Learning and refining Kabaddi techniques requires continuous feedback on
body posture, joint coordination, and temporal execution of movements.
Existing training practices depend heavily on subjective observation by
coaches or self-assessment through video playback, both of which lack
quantitative evaluation. These approaches make it difficult to
consistently identify precise movement errors, especially minor
joint-level deviations that significantly impact technique effectiveness
and balance during gameplay.

There is a clear need for a training system that can objectively analyze
human movement, compare it against expert-defined standards, and present
actionable insights without requiring constant human supervision. Such a
system must handle variations in execution speed, body scale, and
viewing conditions while maintaining consistency in evaluation. The
absence of an accessible, pose-driven feedback mechanism limits
independent practice and slows skill acquisition, highlighting the
necessity for an intelligent, explainable, and structured training
solution tailored specifically for Kabaddi movements.
\end{quote}

\textbf{1.3 OBJECTIVES}

\begin{quote}
The primary objective of this project is to design and develop an
intelligent Kabaddi training system that enables objective evaluation of
player movements using pose-based analysis.

To capture and process user-performed Kabaddi actions in a structured
manner suitable for computational analysis.

To normalize and align human pose data in order to handle variations in
execution speed, body proportions, and camera perspectives.

To quantitatively measure deviations between expert reference movements
and user executions at both joint and temporal levels.

To localize movement errors across different phases of a technique,
enabling precise identification of incorrect posture or timing.

To generate interpretable performance metrics that can support
meaningful LLM feedback and guide technique improvement.
\end{quote}

\textbf{1.4 SOLUTION OVERVIEW}

\begin{quote}
The proposed solution is an augmented reality--assisted training system
that evaluates Kabaddi techniques through structured pose analysis
rather than direct video comparison. The system is designed to separate
visual demonstration, movement capture, and performance evaluation into
well-defined stages, ensuring clarity and consistency in assessment.
Expert movements are represented as reference pose sequences, which
serve as standardized benchmarks for evaluating user-performed actions.

At the input level, user practice sessions are recorded using a mobile
device, and the captured video is treated solely as a source for pose
extraction. Raw video frames are not used directly for comparison or
scoring. Instead, pose estimation outputs are converted into a
normalized skeletal representation, enabling stable analysis independent
of lighting conditions, background variations, or camera placement.

To address differences in execution speed and motion pacing, the
solution incorporates temporal alignment techniques that synchronize
user movements with expert reference sequences. This alignment ensures
that fast or slow executions are evaluated fairly by comparing
corresponding motion phases rather than absolute time frames. As a
result, performance evaluation focuses on movement quality instead of
execution duration.

Following alignment, the system performs joint-wise and frame-wise error
computation to quantify deviations between the user and expert poses.
These errors are aggregated into structured metrics that highlight
spatial inaccuracies and temporal inconsistencies across different
phases of a technique. This approach allows precise localization of
movement faults without relying on subjective interpretation.

The final output of the system consists of quantitative performance
scores and interpretable error indicators that form the basis for
intelligent feedback. By maintaining a strict separation between
evaluation and interpretation, the solution ensures that feedback
generation is explainable, extensible, and grounded in measurable data.
Overall, the proposed approach provides a systematic and objective
framework for skill assessment and improvement in Kabaddi training.
\end{quote}

\textbf{1.5 ORGANIZATION OF THE REPORT}

\begin{quote}
The report is organized into six chapters, each addressing a specific
aspect of the project in a structured and logical manner, supported by
appropriate system diagrams and technical explanations. The chapters are
outlined as follows.

\textbf{CHAPTER 1} introduces the context and motivation behind the
proposed AR-based Kabaddi training system. It presents the problem
statement, clearly defining the limitations of existing training
approaches, and outlines the objectives of the project. This chapter
also provides a high-level overview of the proposed solution and
establishes the scope and direction of the work.

\textbf{CHAPTER 2} presents a comprehensive literature survey of
existing research related to human pose estimation, motion analysis,
temporal alignment techniques, and technology-assisted sports training
systems. It reviews prior methodologies, highlights key observations
drawn from existing work, discusses their limitations, and identifies
the research gap that motivates the proposed approach.
\end{quote}

\textbf{Chapter 2}

\textbf{LITERATURE SURVEY}

\begin{quote}
This chapter presents a focused review of existing research related to
human movement analysis, pose estimation, and intelligent feedback
systems in sports and activity training contexts. The surveyed
literature highlights the progression from traditional video-based
analysis methods to structured pose-driven and learning-based
approaches, emphasizing objective evaluation and quantitative assessment
of movement quality. By examining relevant studies, this chapter
establishes the technical foundation upon which this project is built
and identifies how prior work has addressed similar challenges in motion
comparison and feedback generation. The insights derived from these
studies help position the proposed system within the broader research
landscape while clarifying the specific gaps that motivate this work.
\end{quote}

\textbf{2.1 EXISTING WORK}

\begin{quote}
Human motion analysis has been extensively studied across biomechanics,
activity recognition, and sports performance evaluation, with increasing
emphasis on markerless and pose-based representations. Aleksić et al.
{[}1{]} conducted a detailed validation of markerless pose estimation by
comparing it against gold-standard 3D marker-based motion capture
systems for biomechanical temporal analysis. Their work demonstrates
strong agreement between the two approaches, indicating that skeletal
joint trajectories derived from markerless systems are sufficiently
accurate for quantitative motion analysis. However, the study is
restricted to controlled jump-based tasks and relies on multi-camera
setups, limiting its applicability to unconstrained, sport-specific
movements performed in real-world environments.

Pose estimation combined with machine learning techniques has also been
applied to temporal activity analysis. Lin et al. {[}2{]} proposed a
pipeline integrating OpenPose with recurrent neural network models for
fall detection from video sequences. The study confirms that
pose-derived temporal features can effectively capture human motion
dynamics and achieve high classification accuracy. While this work
validates the robustness of pose-based temporal modeling, it focuses on
event detection rather than detailed evaluation of movement quality or
comparison against expert reference executions, which are essential for
technique-oriented sports training systems.

Skeleton-based action recognition has further advanced the use of
structured skeletal representations. Devzel et al. {[}3{]} introduced a
one-shot skeleton-based action recognition framework using graph
convolutional networks for strength and conditioning exercises. Their
approach demonstrates that skeletal motion patterns can be
discriminative even with limited training data. Despite its
effectiveness in recognizing action categories, the framework emphasizes
classification outcomes and does not address continuous similarity
measurement, temporal alignment, or joint-level error localization
required for movement correction and coaching applications.

Real-time sports performance analysis has been explored by Long Liu et
al. {[}4{]}, who integrated OpenPose with DeepSORT to enable
multi-person pose estimation and tracking in volleyball scenarios. The
system handles occlusion and player interaction while maintaining
real-time performance, supporting the feasibility of pose-based
pipelines in dynamic sports environments. However, the primary focus
remains on tracking and performance statistics rather than
individualized comparison of player movements against expert-defined
standards.

Mixed reality--based instructional systems have also leveraged pose and
motion data to enhance physical training. Ihara et al. {[}5{]} proposed
a framework that converts 2D instructional videos into three-dimensional
mixed reality avatars by extracting motion trajectories and rendering
them in an immersive environment. This approach highlights the
effectiveness of visual augmentation for learning physical skills.
Nevertheless, it requires specialized hardware and does not provide
automated, quantitative evaluation or personalized corrective feedback
based on skeletal similarity metrics.

A critical methodological component across motion comparison studies is
temporal alignment of movement sequences. Human actions are often
performed at varying speeds while preserving their structural order,
making direct frame-to-frame comparison unreliable. Dynamic Time Warping
(DTW) has been established as a standard non-linear temporal alignment
technique for such scenarios. Müller {[}6{]} formally presented DTW as a
robust method for aligning motion and gesture sequences with different
temporal rates, and it has since been widely adopted in motion analysis,
gait studies, and biomechanics. Keogh and Ratanamahatana {[}7{]} further
reinforced DTW's theoretical foundation for exact alignment of
time-series data with temporal distortions.

In biomechanics and pose-based motion analysis, DTW is frequently used
either explicitly or implicitly to establish temporal correspondence
between joint trajectory sequences prior to comparison. Studies such as
Aleksić et al. {[}1{]}, which compare time-series joint data across
different measurement systems, inherently require temporal normalization
to ensure meaningful correspondence. DTW represents the standard
non-linear solution for this requirement in motion analysis literature.
Following temporal alignment, spatial similarity between poses is
commonly quantified using Euclidean distance between corresponding joint
positions or angles. This combination of DTW for temporal alignment and
Euclidean distance for spatial comparison is a well-established and
interpretable strategy in human motion analysis and action comparison
tasks {[}6{]}, {[}8{]}. It enables objective evaluation of movement
quality by separating timing variability from structural pose deviation,
thereby forming a reliable basis for pose-driven training and assessment
systems.
\end{quote}

\textbf{2.2 OBSERVATIONS FROM THE EXISTING WORK}

\begin{quote}
The reviewed literature consistently demonstrates that pose-based
representations are effective for modelling human motion across a wide
range of applications, including biomechanics, activity recognition, and
sports performance analysis. Markerless pose estimation systems have
been shown to provide sufficiently accurate skeletal data for
quantitative analysis when compared against traditional motion capture
setups. This establishes skeletal joint trajectories as a reliable
abstraction for movement evaluation, particularly in scenarios where
intrusive sensors or multi-camera installations are impractical.

Another key observation is the widespread reliance on temporal modelling
to handle variability in human motion execution. Studies involving
action recognition, performance monitoring, and biomechanical comparison
implicitly or explicitly account for differences in execution speed
through temporal normalization or alignment techniques. Dynamic Time
Warping emerges as a recurring methodological choice in classical motion
analysis literature for synchronizing time-series data, enabling fair
comparison of motion phases rather than absolute time indices.

The literature also reveals a strong emphasis on classification and
detection tasks, such as activity recognition, fall detection, or action
labelling. While these approaches successfully identify what action is
being performed, they provide limited insight into how well an action is
executed. Most systems prioritize categorical outcomes or event-level
detection, offering minimal support for continuous similarity scoring or
fine-grained assessment of movement quality.

Furthermore, visualization-oriented systems, including augmented and
mixed reality--based training frameworks, highlight the pedagogical
value of visually demonstrating expert movements. These systems improve
user understanding through immersive representations but often lack a
rigorous quantitative evaluation layer. Feedback is primarily visual or
descriptive, without explicit linkage to measurable joint-level
deviations or temporal inconsistencies.

Overall, the existing work indicates a clear separation between systems
that excel at pose extraction and temporal modelling and those that
provide instructional visualization. There is limited integration of
robust temporal alignment, precise spatial error computation, and
interpretable feedback within a single framework. This gap highlights
the need for training systems that move beyond detection and
visualization toward objective, pose-driven evaluation of movement
quality.
\end{quote}

\textbf{2.3 LIMITATIONS OF EXISTING WORK}

\begin{quote}
Despite significant progress in pose estimation and motion analysis,
several limitations remain evident in existing research. Many pose-based
systems are evaluated under controlled conditions or rely on constrained
datasets, which limits their robustness when applied to unconstrained,
real-world sports movements. Variations in camera placement, execution
style, and environmental factors often degrade performance, making it
difficult to generalize these systems for independent training scenarios
without careful adaptation.

Another notable limitation is the predominant focus on action
recognition and event detection rather than movement quality assessment.
Most existing approaches are designed to classify actions or detect
anomalies, offering binary or categorical outputs. Such formulations do
not support fine-grained evaluation of technique execution, nor do they
provide detailed insights into joint-level deviations or phase-wise
errors that are essential for corrective sports training.

Furthermore, while temporal alignment techniques such as Dynamic Time
Warping are well-established, they are often treated as auxiliary
preprocessing steps rather than being integrated into a structured
evaluation framework. Many systems either assume temporal correspondence
or rely on fixed-length normalization, which can obscure meaningful
motion discrepancies when execution speeds vary significantly between
subjects.

Finally, feedback mechanisms in current systems are largely limited to
visual overlays or summary indicators. There is a lack of explainable,
data-driven feedback that clearly links quantitative error metrics to
specific movement faults. This disconnect between analysis and
interpretation reduces the practical usefulness of existing solutions
for skill refinement, particularly in technique-intensive sports such as
Kabaddi.
\end{quote}

\textbf{2.4 SUMMARY OF LITERATURE SURVEY}

\begin{quote}
The literature survey highlights a clear evolution from traditional
video-based observation methods toward pose-driven and learning-based
approaches for human motion analysis. Prior research confirms that
markerless pose estimation can reliably capture skeletal joint
trajectories suitable for quantitative evaluation, providing a practical
alternative to sensor-based motion capture systems. Temporal modeling
and alignment have been identified as essential components for handling
natural variability in human movement, with Dynamic Time Warping
emerging as a well-established solution for synchronizing motion
sequences executed at different speeds.

At the same time, existing systems predominantly emphasize action
recognition, detection, or high-level performance monitoring. While
these approaches are effective in identifying activities or abnormal
events, they offer limited support for detailed assessment of movement
quality and technique correctness. Visualization-focused training
systems improve interpretability through augmented or mixed reality
representations but often lack a rigorous, quantitative evaluation layer
that links visual feedback to measurable errors.

Overall, the surveyed work reveals a gap between accurate pose
extraction, robust temporal alignment, and interpretable movement
evaluation within a unified framework. There is limited integration of
joint-level error computation, phase-wise analysis, and objective
similarity scoring tailored for technique-oriented sports training.
These observations motivate the need for a structured pose-based
evaluation approach that combines established alignment methods with
precise spatial error metrics, forming a principled foundation for
intelligent and explainable training support systems.
\end{quote}

\end{document}
